\section{Further research questions and proposed next work}
\label{sec:research}

The following questions are ordered by how directly they build on the
proved and implemented results. Each asks for a concrete mathematical
or computational advance, with a criterion by which progress can be assessed.

\subsection{Transport an embedded boundary pattern through a no-fill schedule}

Can the immutable donor structure be lifted from unsigned rows to ordered
strand systems on the boundary of the handlebody? The input should specify
actual source endpoints, local cyclic orders and fibre transport, not only
intersection counts. The output should encode the induced pattern after
each geometric cancellation or directly on the final handlebody.

The passenger theorem bounds the growth of numerical counts. The unresolved
issue is the geometry of their concatenation. A successful result should
give a polynomial bound on the size of a shared representation and an
independent replay theorem. Small counterexamples with equal incidence
and different boundary order should be part of the validation set.

\subsection{Give saturated bands a compressed marked realization}

The saturated theorem omits a potentially enormous family of even band
columns because they cannot be pivots. Can one retain those bands as a
small typed family while transporting specified markings through them?
The target is a source-certified relative band system, with attachment
addresses and cyclic order, whose encoding is polynomial in the source
and the bit length of the base topology.

An abstract free-group presentation with a large product of commutators
does not settle this geometric problem. The representation should support
both orientable and nonorientable bases and arbitrary compatible annular
gluing twists. The punctured-torus and M\"obius examples provide elementary
test cases, while the seven new source cuts provide native regression cases.

\subsection{Identify small repairs of a raw stopping core}

Which geometric moves destroy a specified stopping set by decreasing
actual intersections? The persistence theorem proves that searching more
positive pivot orders cannot help. A useful repair theorem would identify
a bounded-size configuration in every blocked core belonging to a
specified topological class and supply a verified transformation that
reduces a stated potential.

The class should be restricted first. For example, study layered solid
tori or source cuts with a bounded number of saturated relations. The
goal is a theorem with explicit dependence on the chosen parameter,
not a promise inferred from random small examples. The blocked ball
matrix $\left(\begin{smallmatrix}1&1\\1&2\end{smallmatrix}\right)$ is the
first required test: a useful repair can undo a handle slide, but no
positive raw pivot order can finish it.

\subsection{Optimize the choice of source forest with certified scope}

Our completeness theorem fixes the core forest. How does
$\Stop(A)$ change when a tree edge is exchanged with a nonforest edge?
An incremental exact incidence update could avoid reconstructing the
entire source for every forest. The difficult part is to prove the
geometric cocore correspondence and to bound a complete family of
forests or exchanges.

A practical first deliverable is a source-certified one-edge exchange
with an independently verified new aggregate matrix. A later theorem
could characterize a class in which polynomially many exchanges expose
full raw cancellation. The mere existence of a successful forest does
not give a fast complete search over all forests.

\subsection{Use a stopping certificate to focus higher-level moves}

The residual is a canonical row set in a fixed matrix. Can repair proposals
be restricted to neighborhoods of that set while preserving completeness?
One possibility is to keep a typed incidence graph containing its rows,
their columns and the nearby geometric charts. A theorem would need to
show that irrelevant regions can be retained by identity while a useful
repair is found inside the selected region.

This is a locality question about geometric simplification, not a
consequence of the algebraic stopping condition alone. Its benefit would
be a smaller complete move neighborhood with an explicit certificate of
unchanged exterior data.

\subsection{Extend the consumer to three-handle cores}

Ten cuts in the retained audit still have core three-handles. Can a dual
source peeling rule identify geometric two/three-handle cancellations
without flattening the compressed cut? A suitable rule must authenticate
the attaching sphere and its unique intersection, rather than treating
an algebraic incidence as sufficient.

Begin with cases where a genuine free face is present and record its
effect on all interface charts. Then investigate whether a bounded
sequence of dual moves covers a useful larger class. Every accepted
reduction should have an independent source replay and preserve the
component census used by the present implementation.

\subsection{Resolve the eleven remaining raw stalls}

The source audit gives eleven cuts with a failed combined raw system.
Their exact canonical stopping witnesses are included: six obstructions
already lie in the original core, and five require added saturated disk
relations. Classify the geometry behind those witnesses: shared annular
powers, interacting original two-handles, graph cycles, or another
configuration. This is a finite, concrete next experiment.

The desired output is not merely a longer list of heuristic moves. It
is a small set of proved geometric lemmas, each supported by source
certificates, that either resolves a family of stalls or explains why
a new surface or disk system is needed. Some test cuts may actually
fail to be handlebodies; in those cases a separate authenticated
topological obstruction is the appropriate goal. Any cases still
unsupported should remain inconclusive.

\subsection{Unify geometric source compilation with the incoming interface kernels}

The moment-coded first-return report compresses boundary observations,
while the present operator consumes their ordered annular words. Can the
source proofs be reused without changing their trust boundary, reducing
the repeated orbit-query cost that dominates small examples?

The interface contract must preserve shared corners, actual gaps,
orientation transport, whole-circle detection and port ownership. A
useful implementation should retain the old query as a differential
oracle and record both complete-query and finite-kernel timings. This
question is promising precisely because the current empirical evidence
does not show a broad whole-query speedup from the new core planner alone.

\subsection{Find a useful-surface coverage theorem with an explicit budget}

Can every nonterminal member of a restricted hierarchy class be advanced
by a surface whose active interface, saturated residue and repair cost
have a provable bound? Candidate parameters include the number of
saturated relations and the geometry of their stopping core. A useful
theorem should specify how the surface is found, how failures are
charged and how the inherited pattern is transported.

An existence theorem without a producer or an upper bound on failed
candidates is insufficient for the running-time target. Conversely, a
complete producer on a clearly defined restricted class would already
be a substantial result, even before extension to all knot exteriors.

\subsection{Prove a global branching and state-size potential}

Find a potential controlling branching, resets and encoded state size.
Theorem~\ref{thm:global-contract} asks for an $O((\log n)^2)$ path
branching budget with polynomial depth and per-state work.

The potential must include boundary patterns, geometric repairs and
source rebuilding. A raw stopping core alone cannot suffice: coordinate
changes can create or remove it. Proving this combined bound is the
principal remaining step toward a general $n^{O(\log n)}$ recognizer.
